 \chapter{Temps d'attente et chaînes de Markov}\label{ch:iv1:waiting-times-and-markov-chains}

\begin{ivtrains}
Comportement à long terme et temps d'attente dans les processus aléatoires~: combien de temps faut-il attendre avant qu'un événement se produise, lequel de deux événements
survient en premier, et où le processus passe son temps. Les questions suivantes demandent de repérer les bons états, ou une symétrie qui élimine la plus grande partie
de l'algèbre.
\end{ivtrains}

\begin{question}[$\star$ \iqroles{QR, QT, Trader, Analyste} \iqfirm{toutes} \iqlevel{junior}]\label{iq:iv1:waiting-times-and-markov-chains:1}
Vous lancez une pièce équilibrée jusqu'à ce que deux piles apparaissent de suite. Quelle est l'espérance du nombre de lancers~? Comment cela change-t-il si vous attendez plutôt
l'apparition de pile suivi immédiatement de face~?
\end{question}

\begin{question}[$\star$ \iqroles{QR, QT, Trader, Analyste} \iqfirm{toutes} \iqlevel{junior}]\label{iq:iv1:waiting-times-and-markov-chains:2}
Vous lancez un dé équilibré jusqu'à obtenir deux six de suite. Quelle est l'espérance du nombre de lancers~? Et pour trois piles de suite avec une pièce équilibrée~?
\end{question}

\begin{question}[$\star$ \iqroles{QR, QT, Trader, Analyste, Risque} \iqfirm{toutes} \iqlevel{junior}]\label{iq:iv1:waiting-times-and-markov-chains:3}
Vous et un ami lancez tour à tour un dé équilibré, et c'est vous qui commencez. Le premier à obtenir un six gagne. Quelle est votre probabilité de gagner~?
\end{question}

\begin{question}[$\star$ \iqroles{QR, QT, Trader} \iqfirm{toutes} \iqlevel{junior}]\label{iq:iv1:waiting-times-and-markov-chains:4}
Deux joueurs sont à égalité, et le premier à prendre deux points d'avance gagne. Le joueur A gagne chaque point avec la probabilité $0.6$, indépendamment des autres. Quelle est la
probabilité que A gagne, et combien de points supplémentaires sont joués en moyenne~?
\end{question}

\begin{question}[$\star\star$ \iqroles{QR, QT, DS, Risque} \iqfirm{toutes} \iqlevel{junior}]\label{iq:iv1:waiting-times-and-markov-chains:5}
Un processus possède deux états. Dans l'état A, il reste avec la probabilité $0.9$ à chaque pas et passe sinon dans B. Dans l'état B, il reste avec
la probabilité $0.7$, et passe sinon dans A. Dans le long terme, quelle fraction du temps est-on dans B~? Quelle est la durée moyenne d'un séjour dans B~?
Une fois le processus installé dans son comportement à long terme, quelle est la corrélation entre «~être dans B maintenant~» et «~être dans B après $n$ pas~»~?
\end{question}

\begin{question}[$\star\star$ \iqroles{QR, QT, Trader} \iqfirm{toutes} \iqlevel{junior}]\label{iq:iv1:waiting-times-and-markov-chains:6}
Une pièce équilibrée est lancée jusqu'à ce que l'une des séquences pile-pile-face ou face-pile-pile apparaisse sur trois lancers consécutifs. Vous gagnez si pile-pile-face sort
en premier. Quelle est votre probabilité de gagner~?
\end{question}

\begin{question}[$\star\star$ \iqroles{QR, QT, Trader} \iqfirm{teneur de marché, hedge fund, banque} \iqlevel{junior, confirmé}]\label{iq:iv1:waiting-times-and-markov-chains:7}
Vous lancez une pièce équilibrée jusqu'à ce que le motif PFPFP apparaisse. Quelle est l'espérance du nombre de lancers~? Quelle est-elle pour PPPFF, et pourquoi les deux valeurs diffèrent-elles alors que chaque motif a la même chance d'apparaître à une position donnée~?
\end{question}

\begin{question}[$\star\star$ \iqroles{QR, QT, Trader, QD} \iqfirm{toutes} \iqlevel{junior}]\label{iq:iv1:waiting-times-and-markov-chains:8}
Une fourmi part d'un sommet d'un cube. À chaque pas, elle parcourt au hasard l'une des trois arêtes issues de ce sommet. En moyenne, combien de
pas met-elle pour atteindre le sommet diagonalement opposé~?
\end{question}

\begin{question}[$\star\star$ \iqroles{QR, QT, Trader} \iqfirm{hedge fund, teneur de marché} \iqlevel{junior}]\label{iq:iv1:waiting-times-and-markov-chains:9}
Vous lancez un dé équilibré à répétition et tenez un total cumulé. En moyenne, combien de lancers faut-il pour que le total soit divisible par 6~?
\end{question}

\begin{question}[$\star\star$ \iqroles{QR, QT, Trader} \iqfirm{banque, hedge fund} \iqlevel{junior}]\label{iq:iv1:waiting-times-and-markov-chains:10}
Vous lancez un dé équilibré à répétition et tenez un total cumulé. En moyenne, combien de lancers faut-il pour que le total atteigne 10 ou plus~? Donnez une méthode exacte et une approximation rapide
pour la vérifier, puis celle pour un total de 100 ou plus.
\end{question}

\begin{question}[$\star\star$ \iqroles{QR, QT, Trader} \iqfirm{toutes} \iqlevel{junior}]\label{iq:iv1:waiting-times-and-markov-chains:11}
Vous lancez à répétition une paire de dés équilibrés et observez leur somme. Quelle est la probabilité qu'une somme de 12 apparaisse avant deux sommes consécutives égales à 7~?
\end{question}

\begin{question}[$\star\star$ \iqroles{QR, QT, Risque, DS, Trader} \iqfirm{toutes} \iqlevel{junior}]\label{iq:iv1:waiting-times-and-markov-chains:12}
Les bus quittent un arrêt les uns après les autres, et chaque intervalle entre deux bus vaut indépendamment 5 ou 15 minutes, avec la même probabilité. Vous arrivez à un instant aléatoire. Quelle est votre attente moyenne pour le bus suivant~?
\end{question}

\begin{question}[$\star\star$ \iqroles{QR, QT, Risque} \iqfirm{toutes} \iqlevel{junior, confirmé}]\label{iq:iv1:waiting-times-and-markov-chains:13}
Une bactérie meurt avec la probabilité $1/4$, survit sans changement avec la probabilité $1/4$, ou se divise en deux avec la probabilité $1/2$. Chaque bactérie procède ainsi indépendamment, une fois par minute. En partant d'une seule bactérie, quelle est la probabilité que la colonie finisse par disparaître~?
\end{question}

\begin{question}[$\star\star\star$ \iqroles{QR, QT, Trader} \iqfirm{toutes} \iqlevel{junior, confirmé}]\label{iq:iv1:waiting-times-and-markov-chains:14}
Vous continuez à lancer un dé équilibré et à ajouter les résultats jusqu'à ce que votre total atteigne 1\,000 ou plus. Quel total final est le plus probable, et quelle est sa probabilité, à trois décimales près~?
\end{question}

\begin{question}[$\star\star\star$ \iqroles{QR, QT, Trader} \iqfirm{teneur de marché, hedge fund, crypto} \iqlevel{junior, confirmé}]\label{iq:iv1:waiting-times-and-markov-chains:15}
Vous et un adversaire nommez chacun un motif de trois lancers de pièce, tel que PFP, et votre adversaire nomme le sien en premier. On lance une pièce équilibrée jusqu'à ce que l'un des deux motifs apparaisse, et celui qui a nommé ce motif gagne. Si votre adversaire nomme PPP, que devriez-vous nommer, et quelle est votre chance de gagner~? Et s'il nomme PFP~? Existe-t-il une réponse qui vous garantit une chance fixe minimale, quel que soit le motif nommé~?
\end{question}

\begin{question}[$\star\star\star$ \iqroles{QR, QT, Trader} \iqfirm{teneur de marché, hedge fund} \iqlevel{junior, confirmé}]\label{iq:iv1:waiting-times-and-markov-chains:16}
Vous vous trouvez sur le chiffre 12 du cadran d'une horloge. À chaque pas, vous passez à l'un des deux chiffres voisins avec la même probabilité, et vous continuez jusqu'à avoir visité tous les chiffres. Quelle est la probabilité que le 6 soit le dernier nouveau chiffre visité~? En moyenne, combien de pas dure la marche complète~?
\end{question}

\begin{question}[$\star\star\star$ \iqroles{QR, QT, Trader} \iqfirm{hedge fund, teneur de marché, banque} \iqlevel{junior, confirmé}]\label{iq:iv1:waiting-times-and-markov-chains:17}
Un cavalier part d'un coin d'un échiquier vide et se déplace au hasard, en choisissant à chaque fois avec la même probabilité parmi ses coups légaux. En moyenne, combien de coups faut-il pour revenir dans ce coin~?
\end{question}

\begin{question}[$\star\star\star$ \iqroles{QR, QT} \iqfirm{hedge fund, teneur de marché} \iqlevel{confirmé}]\label{iq:iv1:waiting-times-and-markov-chains:18}
Un jeton se trouve sur un coin d'une grille de 3 par 3 cases reliées par leurs côtés, soit 9 positions et 12 liaisons entre cases voisines. À chaque pas, il passe à une position voisine choisie uniformément. En moyenne, combien de pas faut-il au jeton pour atteindre le coin opposé~?
\end{question}

\begin{question}[$\star\star\star$ \iqroles{QR, QT} \iqfirm{hedge fund, teneur de marché} \iqlevel{junior, confirmé}]\label{iq:iv1:waiting-times-and-markov-chains:19}
Vous lancez une pièce équilibrée jusqu'à ce que deux piles apparaissent de suite. L'espérance du nombre de lancers vaut 6. Quelle est la variance du nombre de lancers~? Comment cela se compare-t-il à l'attente de pile suivi de face~?
\end{question}

\begin{question}[$\star\star\star$ \iqroles{QR, QT} \iqfirm{hedge fund, teneur de marché} \iqlevel{confirmé}]\label{iq:iv1:waiting-times-and-markov-chains:20}
Vous mélangez un jeu de 52 cartes en prenant à répétition la carte du dessus et en la réinsérant à l'une des 52 positions possibles, choisie uniformément, y compris celles du dessus et du dessous. Suivez la carte qui se trouvait tout en bas au départ. En moyenne, combien de mouvements faut-il pour que cette carte atteigne le dessus puis soit réinsérée~? Pourquoi le jeu est-il alors parfaitement mélangé~?
\end{question}

\begin{question}[$\star\star\star$ \iqroles{QR, QT} \iqfirm{hedge fund, teneur de marché} \iqlevel{confirmé}]\label{iq:iv1:waiting-times-and-markov-chains:21}
Dix pièces reposent sur une table, toutes sur face. À chaque pas, vous choisissez uniformément une pièce au hasard et la retournez. En moyenne, combien de pas faut-il pour que toutes les pièces soient sur pile~? Et combien de pas faut-il pour revenir pour la première fois à toutes les pièces sur face~?
\end{question}

\begin{question}[$\star\star\star\star$ \iqroles{QR, QT} \iqfirm{teneur de marché, hedge fund} \iqlevel{confirmé}]\label{iq:iv1:waiting-times-and-markov-chains:22}
Vous et un adversaire nommez chacun un motif de quatre lancers de pièce, et votre adversaire nomme le sien en premier. On lance une pièce équilibrée jusqu'à ce que l'un des motifs apparaisse, et son propriétaire gagne. Votre adversaire nomme PFPP. Quel motif vous donne la meilleure chance, quelle est cette chance, et quelle est l'espérance du nombre de lancers avant la fin de la partie~?
\end{question}

\begin{question}[$\star\star\star\star$ \iqroles{QR, QT} \iqfirm{teneur de marché, hedge fund} \iqlevel{confirmé}]\label{iq:iv1:waiting-times-and-markov-chains:23}
Vous lancez un dé équilibré et vous arrêtez dès que les six derniers lancers sont tous différents, c'est-à-dire lorsqu'ils contiennent chaque face une fois. Quelle est l'espérance du nombre de lancers~?
\end{question}

\begin{question}[$\star\star\star\star$ \iqroles{QR, QT, Risque} \iqfirm{hedge fund, crypto} \iqlevel{confirmé}]\label{iq:iv1:waiting-times-and-markov-chains:24}
Une cellule meurt avec la probabilité $1/3$, ou se divise en deux cellules avec la probabilité $2/3$. Chaque cellule se comporte indépendamment. Vous partez d'une seule cellule. Quelle est la probabilité que la famille finisse par disparaître~? Sachant qu'elle disparaît, quelle est l'espérance du nombre total de cellules ayant jamais vécu, en comptant la première~?
\end{question}

\begin{question}[$\star\star\star\star$ \iqroles{QR, QT, Risque} \iqfirm{hedge fund} \iqlevel{confirmé}]\label{iq:iv1:waiting-times-and-markov-chains:25}
Chaque individu d'une population a indépendamment $k$ enfants avec la probabilité $(1/2)^{k+1}$, pour $k = 0, 1, 2, \dots$, puis meurt, si bien qu'un enfant est la moyenne. Vous partez d'un seul individu à la génération 0. Quelle est la probabilité que la génération 10 soit vide~? Sachant qu'elle ne l'est pas, quelle est son espérance de taille~? Quelle est l'espérance de la génération à laquelle la famille disparaît~?
\end{question}
